\chapter{Conclusão}
\label{cap:conclusao}

Este trabalho replicou o FactorVAE --- modelo de fatores latentes probabilístico baseado em inferência variacional --- no mercado brasileiro de ações e avaliou seu desempenho em termos de \textit{skill} cross-sectional, performance de portfólio ajustada ao risco e confronto com os resultados de \textcite{duan2022} no mercado A-share chinês. A pergunta que orientou a investigação é se o FactorVAE mantém vantagem de \textit{skill} sobre os baselines na B3 e se essa vantagem sobrevive à tradução em portfólio sob custos de transação realistas.

O FactorVAE obteve Rank IC médio de 0{,}040 no período de teste, contra 0{,}035 do Conditional Autoencoder, 0{,}025 do GRU e 0{,}016 do IPCA. A vantagem sobre o GRU --- o baseline neural mais próximo em complexidade --- é de 60\% em termos relativos, e o Rank ICIR (0{,}246 vs.\ 0{,}217) confirma que essa vantagem se sustenta ao longo dos sete anos de teste. A combinação de estrutura fatorial explícita e inferência variacional, implementada pelo esquema professor-aluno, produz um sinal de ordenação cross-sectional mais forte e mais consistente do que o de uma GRU pura treinada com a mesma entrada. Na estratégia de portfólio TopK-Drop ($k = 50$, custo de 25~bps \textit{one-way}), o FactorVAE liderou entre os modelos avaliados em CAGR (9{,}97\%), registrou a menor volatilidade (19{,}72\%) e o menor drawdown máximo (41{,}24\%). O Sharpe ficou em $-0{,}002$ com taxa de referência de 10\% a.a., e o excesso anualizado sobre o Ibovespa foi de $-0{,}81$~p.p. --- o que configura ausência de alfa líquido contra o benchmark após custos. A vantagem do FactorVAE sobre os demais modelos se materializa, portanto, na dimensão de risco e não na de retorno absoluto.

O confronto com \textcite{duan2022} revela convergência qualitativa e divergência em magnitudes. A ordenação dos modelos se reproduz: o FactorVAE lidera em Rank IC sobre os demais modelos de fatores e sobre os baselines neurais em ambos os estudos. Os valores absolutos, porém, são menores na B3 (Rank IC de 0{,}040 vs.\ 0{,}055; ICIR de 0{,}246 vs.\ 0{,}568; alfa próximo de zero vs.\ positivo). Três hipóteses explicam essa diferença. O tamanho do universo é a mais direta: com 300--500 ativos na B3 contra mais de 3.000 no mercado chinês, $k = 50$ representa 10--17\% do cross-section brasileiro vs.\ $\approx$1{,}5\% do chinês, o que dilui mecanicamente o poder de discriminação do sinal. A concentração setorial do mercado brasileiro em commodities, energia e financeiro reduz a dimensionalidade efetiva do cross-section e pode comprimir a dispersão de retornos explorável pelos fatores latentes. O volume de dados de treinamento é menor em ambas as dimensões --- menos ativos por data e período mais curto ---, o que limita a capacidade do modelo de aprender a distribuição condicional dos fatores.

A resposta à pergunta de pesquisa é, portanto, parcial. O FactorVAE mantém vantagem de \textit{skill} sobre todos os baselines avaliados na B3, e essa vantagem se traduz em melhor perfil de risco da estratégia de portfólio. A tradução em alfa líquido positivo após custos, porém, não se concretiza no período observado.

O trabalho possui limitações que condicionam a interpretação dos resultados. As características de entrada são exclusivamente técnicas; a ausência de variáveis fundamentalistas pode limitar o sinal em horizontes nos quais informação contábil é relevante. O custo de transação fixo de 25~bps não captura o impacto de mercado em ativos de menor capitalização, onde os custos efetivos provavelmente excedem essa estimativa. A estratégia \textit{long-only} impede a extração do sinal nos dois extremos do cross-section e limita a avaliação do alfa desalavancado.

Três extensões se conectam diretamente aos resultados obtidos. A incerteza preditiva $\sigma_y^{(i)}$, entregue em forma fechada pelo decodificador, é uma saída da arquitetura que não foi explorada na alocação além da menção ao TDrisk; incorporá-la na função de alocação permitiria avaliar se a representação probabilística dos fatores acrescenta valor econômico além do sinal de retorno esperado. A aplicação do mesmo protocolo a outros mercados emergentes (México, Índia, África do Sul) mapearia como a vantagem do FactorVAE depende de profundidade do universo e concentração setorial, separando efeitos específicos da B3 de limitações inerentes a mercados com cross-section restrito. Por fim, a reavaliação em períodos de teste mais longos, à medida que dados se acumulam, permitiria verificar se a vantagem em Rank IC se mantém fora dos regimes de alta volatilidade que dominam o período amostral.